\documentclass[11pt,letterpaper]{article}

\usepackage[margin=1in]{geometry}
\usepackage{amsmath,amssymb,amsthm}
\usepackage{booktabs}
\usepackage{graphicx}
\usepackage{xcolor}
\usepackage[hidelinks]{hyperref}
\usepackage{microtype}
\usepackage{enumitem}
\usepackage{titlesec}
\usepackage{caption}

\definecolor{mDarkTeal}{HTML}{1c3d6b}
\definecolor{mLightBrown}{HTML}{2f6fb3}

\titleformat{\section}{\large\bfseries\color{mDarkTeal}}{\thesection}{0.6em}{}
\titleformat{\subsection}{\normalsize\bfseries}{\thesubsection}{0.5em}{}
\titlespacing*{\section}{0pt}{1.4em}{0.6em}
\titlespacing*{\subsection}{0pt}{1.0em}{0.4em}

\hypersetup{colorlinks=true, linkcolor=mDarkTeal, citecolor=mDarkTeal, urlcolor=mLightBrown}

\newcommand{\R}{\mathbb{R}}
\newcommand{\Gop}{\mathcal{G}}
\newcommand{\sgrad}{\operatorname{sg}}
\newcommand{\Half}{\tfrac{1}{2}}

\title{\textbf{DNO Surrogate: Diagnostics and Approaches Under Investigation}}
\author{Jonathan Ma}
\date{5 June 2026}

\begin{document}
\maketitle

\begin{abstract}
\noindent
This document is the single working reference for the effort to improve
long-horizon rollout accuracy of a neural surrogate for the 1D Zakharov
water-wave Dirichlet--Neumann operator. It covers: (i) what is wrong with
current rollouts, established by forensic analysis on the existing FNO and
SpectralDNO checkpoints, plus a physical/numerical sensitivity sweep that
rules out integrator and amplitude artifacts; (ii) why the training data
is not the bottleneck; (iii) the three approaches currently in flight ---
a one-step pushforward training loss, a relative-Hamiltonian conservation
loss, and a Craig--Sulem-shaped DNO architecture; and (iv) three
candidates queued for the next iteration.
\end{abstract}

\vspace{0.6em}
\noindent\textbf{Current state (5 June 2026).}
The diagnostic case is closed: the rollout discrepancy is model-driven,
not a data, integrator-resolution, or amplitude artifact. The
training-loss approaches (pushforward, H-cons) reduce error on
soliton-like and modal-instability regimes but regress on broadband
random-sea ICs; the gap is qualitatively similar in both, suggesting
the SpectralDNO architecture is the binding constraint. The
Craig--Sulem-shaped architecture is currently training from scratch
as the load-bearing test of that hypothesis; an H-cons SpectralDNO
is retraining in parallel for an apples-to-apples comparison.

% ===================================================================
\section{Problem statement}

The 1D Zakharov system couples surface elevation $\eta(x,t)$ and surface
velocity potential $\xi(x,t)$ on a periodic domain of length $L$ via the
Dirichlet--Neumann operator $\Gop(\eta;h)$:
\begin{align}
\partial_t \eta &= \Gop(\eta; h)\,\xi,\\
\partial_t \xi &= -g\,\eta - \Half\xi_x^{\,2}
   + \Half\frac{\big(\Gop(\eta)\xi + \eta_x\xi_x\big)^{2}}{1 + \eta_x^{\,2}}.
\end{align}
A neural surrogate $\Gop_\theta$ is trained to predict $\Gop(\eta;h)\xi$
from snapshots of $(\eta,\xi,h)$ and then iterated by an explicit time
integrator at evaluation time. Pointwise prediction is accurate;
long-horizon rollouts are not. The shared baseline objective is the
data-distribution Sobolev-$H^{k}$ loss
\[
  \mathcal{L}_{0} = \big\| \Gop_\theta(\eta;h)\xi - \Gop_{\mathrm{true}}(\eta;h)\xi \big\|_{H^{k}}^{2},
\]
with $k=1$ in the current configuration. Two surrogates have been trained
on this baseline: a 6-block FNO1d (width 128, modes 64, FiLM depth
conditioning, $\sim 13.1$M parameters) and a 6-block SpectralDNO
(width 128, modes 64, latent 64, $\sim 17.0$M parameters). The SpectralDNO
outputs a residual of the form $C(\eta,h)^{\top} F(\eta,h)\,C(\eta,h)\,\xi$
added to a closed-form linear-DNO baseline, and is therefore linear in
$\xi$ by construction.

% ===================================================================
\section{Diagnostic findings}
\label{sec:diagnostics}

\subsection{Failure mode: coherent in-band error}

Across the worst finite-error initial conditions in the 10-regime eval
suite (Tanaka, Benjamin--Feir, Stokes, random sea, linear), rel-$L^{2}$
error grows linearly or with slight super-linearity in time
($r^{2}_{\mathrm{lin}}$ between 0.53 and 0.93). The error spectrum is
dominated by the same wavenumber band as the truth spectrum (per-IC
cosine similarity 0.57--1.00 for FNO); the high-$k$ ``Gibbs'' fraction is
below 0.18 in every worst case. NaN blow-ups in bf\_g1 and bf\_modal
occur at $t = 81$--$189$ of $T_{\max}=200$ while $|\eta|_{\max}$ stays
small (0.04--0.12). This is \emph{numerical divergence after extended
coherent drift, not a steepening cascade or a high-$k$ artifact}.

A specific smoking gun on the linear regime: FNO at $t=2$ has $99\%$ of
its error in sideband modes $k=10$--$50$, while the truth has all its
energy at $k=6$. The FNO is hallucinating nonlinear sidebands in a
regime where the dynamics are strictly per-mode dispersive --- a
structural consequence of mixing $\xi$ into the network's nonlinear
features. SpectralDNO, which is linear in $\xi$ by construction, cannot
produce this kind of error and dominates FNO on the linear regime by
$-71\%$ in median rel-$L^{2}$.

\subsection{Physics descriptors that predict failure}

Within each regime, the strongest Spearman correlations with final-time
rel-$L^{2}$ are:
\begin{itemize}[leftmargin=1.3em,itemsep=2pt]
  \item Tanaka g0/g1: spectral bandwidth $\rho = +0.76\,/\,+0.78$ (FNO),
        $+0.84\,/\,+0.85$ (DNO+PF).
  \item Random sea (deep): max steepness $\rho = +0.63\,/\,+0.84$.
  \item Stokes (finite): max steepness $\rho = +0.70\,/\,+0.83$.
  \item BF g0/g1: peak wavenumber $k_{\mathrm{peak}}$,
        $\rho \in [+0.33, +0.87]$.
\end{itemize}
Global correlations with $kh$ or depth ($\rho \approx -0.4$) appear strong
but are confounded --- Tanaka has both shallow water and broad spectra.
Within-regime, the load-bearing predictors are \emph{structural}
(bandwidth, steepness, peak wavenumber), not distributional. This is what
high-steepness and broadband ICs need from the model's
nonlinear-in-$\eta$ part --- exactly the part of the current architectures
that is under-parameterised relative to the polynomial structure of the
true DNO.

\subsection{Rollout error vs.\ initial wavenumber}

A per-regime diagnostic picks three ICs per regime binned by initial
mean wavenumber $\langle k\rangle(0) = \sum_k k\,|\hat\eta(k,0)|^{2}/\sum_k
|\hat\eta(k,0)|^{2}$ (low / mid / high) and compares two stacked panels
over the rollout horizon: rel-$L^{2}$ $\eta$-error (log scale, top)
and truth mean wavenumber $\langle k\rangle(t)$ (bottom). On Tanaka g0
the wave's mean wavenumber is essentially stationary in time, and the
rollout error grows two orders of magnitude --- more rapidly for
higher-$k$ ICs. The same pattern holds on the linear regime, where
$\langle k\rangle(t)$ is exactly flat by construction, so the error
growth is not riding on any spectral drift in the underlying wave.
The full panel grid for all 10 regimes is on the companion site.

\subsection{Data audit: not the bottleneck}

A 1000-snapshot sample of the 73.7\,GB training dataset found no
NaN/inf in over $10^{6}$ sampled voxels of $\eta$, $\xi$, or $\Gop\xi$.
Every test regime's IC distribution falls inside the training central
90\% at the median for amplitude, steepness, and dominant wavenumber.
Two minor tails: the maximum-amplitude Tanaka cases nip the training
$p99$ in $\eta$ amplitude; the finite-depth Stokes regime's dominant
wavenumber sits at the upper edge of training $k$ ($\sim 26$, within
the spectral cutoff of $64$). One eval-side issue surfaced: one of the
16 Benjamin--Feir g1 ICs has the \emph{truth} integrator diverge to
NaN at $t \approx 146$ of $T_{\max} = 200$.

\textbf{Conclusion.} The rollout-accuracy ceiling is not a data-coverage
problem. Improvements have to come from the model or the training-time
loss, not from acquiring more data.

\subsection{Linear-regime truth fix}

The legacy ``linear'' regime ICs are physically nonlinear shallow-water
cases ($\eta/h$ up to $0.30$, $kh \approx 0.6$) inherited from a MATLAB
test set. The Zakharov truth solver was diverging on some of them. The
eval suite was updated to use analytic per-mode Fourier propagation of
the linear Zakharov system,
\[
\omega^{2}(k) = g\,k\,\tanh(hk),\qquad
\hat\eta(k,t) = \hat\eta_{0}(k)\cos\omega t + \tfrac{\Gop_{\mathrm{lin}}(k)}{\omega}\,\hat\xi_{0}(k)\sin\omega t,
\]
and analogously for $\hat\xi(t)$, as the linear-regime ``truth''. Amplitude
and top-$k$ structure are preserved to better than $1\%$ over the full
horizon. All linear-regime comparisons in this document use the analytic
truth.

\subsection{Sensitivity sweep: is the FNO discrepancy a numerical artifact?}
\label{sec:sensitivity}

To rule out the obvious physical and numerical explanations for the
rollout discrepancy, we ran a 3-axis sweep on Tanaka g0 case~0
(depth $h{=}0.139$, horizon $T{=}100$, output $\mathrm{d}t{=}1$,
FNO baseline). On each axis we hold the others fixed; the rel-$L^{2}$
$\eta$-error at $T=100$ is:
\begin{center}
\small
\begin{tabular}{lrrrrr}
\toprule
\textbf{Axis} & \multicolumn{5}{l}{\textbf{Values}} \\
\midrule
IC amplitude $\alpha$ & $0.50$ & $0.75$ & $1.00$ & $1.50$ & $2.00$ \\
$\quad$ rel-$L^{2}$ $\eta$ @ $T$ & $0.0059$ & $0.0050$ & $0.0043$ & $0.0080$ & $0.0139$ \\
\midrule
Integrator substeps / output dt & $1$ & $2$ & $4$ & $8$ & $16$ \\
$\quad$ rel-$L^{2}$ $\eta$ @ $T$ & NaN & NaN & $0.0044$ & $0.0043$ & $0.0042$ \\
\midrule
Spatial $N_x$ & $512$ & $1024$ & $2048$ & & \\
$\quad$ rel-$L^{2}$ $\eta$ @ $T$ & $0.1435$ & $0.0043$ & NaN @ $t{=}32$ & & \\
\bottomrule
\end{tabular}
\end{center}
\noindent
\textbf{Read.} The error is roughly $\alpha^{2}$ on the large-amplitude
side (consistent with nonlinear-in-$\eta$ compounding) but stays within
a factor of three across a $4\times$ amplitude range, with a minimum at
the training amplitude $\alpha = 1$. Refining the integrator past
substeps${=}4$ does not change the error to four figures
(the NaN at substeps${\in}\{1,2\}$ is the GL2 stepper itself being
unstable at large dt, unrelated to the FNO). The FNO is, however,
tightly resolution-locked to its training grid: at $N_x{=}512$ it is
$33\times$ worse, and at $N_x{=}2048$ it blows up by $t{=}32$.

\textbf{Conclusion.} The discrepancy is invariant to numerical dt once
the integrator is converged, and only modestly sensitive to physical
amplitude. It is not a numerical artifact and is not amplitude-fragile;
it is a model error. The $N_x$ result is a separate finding ---
SpectralDNO/FNO are nominally resolution-invariant but in practice are
tied to the training grid; this is worth flagging for any future
multi-resolution evaluation.

% ===================================================================
\section{Approach 1: One-step pushforward (training-loss)}

\subsection{Formulation}
Given a data pair $(\eta,\xi)$, advance the surrogate's prediction
through $k$ Euler steps of the Zakharov system, then re-predict $\Gop$ at
the rolled-out state and penalise against the analytical Craig--Sulem
expansion at order four:
\[
  \mathcal{L}_{\mathrm{pf}}
  \;=\; \mathcal{L}\!\Big(
     \Gop_\theta\!\big(\sgrad(\eta_{k}), \sgrad(\xi_{k})\big),\;
     \Gop_{\mathrm{CS}}(\eta_{k}, \xi_{k})
  \Big),
\]
where $(\eta_{j},\xi_{j})$ is the surrogate's $j$-step Euler rollout. The
total loss is $\mathcal{L}_{0} + \lambda_{\mathrm{pf}}\mathcal{L}_{\mathrm{pf}}$.
Gradients flow only through the final prediction at the rolled-out state;
the trajectory $(\eta_{k},\xi_{k})$ and the analytical target are both
stop-gradient-wrapped.

\subsection{Inductive bias}
The single-step baseline only ever exposes the surrogate to
data-distribution inputs. At rollout, the surrogate consumes its own
predictions, where small per-step biases compound. The pushforward loss
forces the surrogate to remain accurate at states it itself produces
under iteration, breaking the covariate-shift compounding.

\subsection{Implementation}
One Euler step of the Zakharov system per pushforward step. The
$\partial_{t}\eta$ term uses the surrogate's own $\Gop$-prediction; the
$\partial_{t}\xi$ term uses the closed-form Zakharov RHS. The analytical
target $\Gop_{\mathrm{CS}}$ is the order-4 Craig--Sulem series with
zero-padded FFTs at pad factor 4. The depth saturation $h_{\max} = 5$
used inside the model is mirrored when building the analytical target so
the target matches the operator the model is approximating.

\subsection{Cost and tunables}
$(1+k)$ model forwards plus one vmapped order-4 Craig--Sulem evaluation
per training step. Active configuration: $k=1$, $\lambda_{\mathrm{pf}}=1$,
$\mathrm{d}t=0.01$.

\subsection{Status}
A 40-epoch run from scratch is complete (best val $0.00490$ at epoch
$31$). Per-regime rollout from that checkpoint: substantial improvement
on Tanaka, Benjamin--Feir, modal, and the linear regime (median rel-$L^{2}$
reductions of 25--70\%) but \emph{regression} on random-sea regimes
(near-zero error at FNO baseline grows to 0.09--0.12). The pushforward
step appears to trade noise-robustness on broadband ICs for fidelity on
coherent soliton-like ICs.

A 40-epoch resume to test for convergence was contaminated by a
checkpoint-restore bug: the saved Adam moments and cosine-LR step
counter did not deserialise cleanly, so the resume re-warmed Adam from
zero and restarted the LR schedule from base lr. The restore path has
since been fixed; an apples-to-apples 80-epoch run can now be done
cleanly when GPU time frees up.

% ===================================================================
\section{Approach 2: Hamiltonian-conservation loss (training-loss)}

\subsection{Formulation}
The Zakharov system is Hamiltonian with
\[
  H(\eta,\xi) = \Half\!\int_\Omega \big( \xi\,\Gop(\eta)\xi + g\,\eta^{2}\big)\,dx.
\]
An exact solution conserves $H$. The loss penalises relative drift of $H$
over one Euler step of the surrogate's own dynamics:
\[
  \mathcal{L}_{H} \;=\; \lambda_{H}\,w(s)\,
  \mathbb{E}\!\left[\min\!\left(\frac{|H_{1}-H_{0}|}{|H_{0}|+\varepsilon},\;C\right)\right],
\]
with $H_{0}$ at the data state, $H_{1}$ at one Euler step (both
$\Gop$-values from the surrogate), linear warmup
$w(s)=\min(s/s_{\mathrm{warm}},1)$, and per-sample clip $C$. Total loss
is $\mathcal{L}_{0} + \mathcal{L}_{H}$.

\subsection{Inductive bias}
Even when pointwise $\Gop$-error per step is small, accumulated rollout
can violate energy conservation, producing non-physical growing or
decaying waves. Penalising $|\Delta H|/|H|$ biases the surrogate toward
predictions consistent with the symplectic structure across one step.
This is a soft physical constraint, not a structural one --- the
architecture is unchanged.

\subsection{Implementation}
Two model forwards (at $(\eta,\xi)$ and at $(\eta_{1},\xi_{1})$), two FFT
pairs (to compute $\xi\Gop(\eta)\xi$ via the spectral inner product),
and one closed-form Zakharov RHS evaluation. The advanced state is
stop-gradient-wrapped; gradients flow only through the two
$\Gop$-predictions. The integrator and the depth saturation are treated
as fixed.

\subsection{Cost and tunables}
Roughly $2\times$ the per-step cost of the single-step baseline. Active
configuration: $\lambda_{H}=1$, $\mathrm{d}t=0.01$,
$s_{\mathrm{warm}}=10\,000$, $C=10$.

\subsection{Status}
A first Modal training run reached epoch 10 (best val 0.00464, ahead
of DNO+PF at the same epoch) before worker preemption. Partial
evaluation on three regimes showed the rollout numbers tracking
DNO+PF (random-sea regression of similar magnitude, slightly worse on
Tanaka median, marginally higher divergence rate) --- no breakthrough
at that early checkpoint. A clean 25-epoch run from scratch is now
training locally for a fair compute-matched comparison.

% ===================================================================
\section{Approach 3: Craig--Sulem-shaped DNO architecture (model)}

\subsection{Formulation}
The true Dirichlet--Neumann operator admits a Craig--Sulem expansion
\[
  \Gop(\eta;h)\xi \;=\; G_{0}(h)\,\xi + G_{1}(\eta,h)\,\xi
                   + G_{2}(\eta\otimes\eta,h)\,\xi + \cdots
\]
where $G_{0}(h) = |D|\tanh(h|D|)$ is the flat-surface linear DNO and each
$G_{n}$ is multilinear in $\eta^{\otimes n}$ and linear in $\xi$, formed
by alternating pointwise multiplications by powers of $\eta$ with
Fourier multipliers $M(D,h)$. The proposed architecture is a sum of
learned Craig--Sulem-shaped blocks:
\[
  \Gop_\theta(\eta;h)\xi \;=\; G_{0}(h)\,\xi
   + \sum_{\ell=1}^{N_{\mathrm{blocks}}}
      \sum_{i=1}^{n_{\mathrm{branches}}}
        M^{\ell,i}_{\mathrm{out}}(D,h)\,\Big[
          \Phi^{\ell,i}(\eta)(x)\;\cdot\;
          \bigl(M^{\ell,i}_{\xi}(D,h)\,\xi\bigr)(x)
        \Big],
\]
where:
\begin{itemize}[leftmargin=1.3em,itemsep=2pt]
  \item $\Phi^{\ell,i}(\eta)$ is a learned pointwise projection from a
        shared $\eta$-feature stack (pointwise powers
        $\eta,\eta^{2},\eta^{3}$ plus spectral derivatives
        $\partial_{x}\eta,\partial_{x}^{2}\eta,|D|^{1/2}\eta,\mathcal{H}\eta$)
        to one spatial channel.
  \item $M^{\ell,i}_{\xi}$ and $M^{\ell,i}_{\mathrm{out}}$ are learned
        depth-and-$k$-aware Fourier multipliers, each parameterised as a
        small MLP of the physical features
        $(k,\,h,\,\tanh(hk),\,k\tanh(hk))$ --- the natural features of
        the linear dispersion relation, so the MLP only learns
        perturbations on the right ansatz.
  \item The $\Phi^{\ell,i}\to$ spatial channel projection is
        zero-initialised so each block starts inert; the model initially
        predicts the linear baseline exactly.
\end{itemize}

\subsection{Inductive bias}
Two simultaneous structural priors:
\begin{enumerate}[leftmargin=1.3em,itemsep=2pt]
  \item \textbf{Linear in $\xi$ by construction.} The output residual is
        a sum of terms each of the form
        $(\text{$\eta$-only})\cdot(M(D,h)\xi)$. This prevents the
        sideband hallucination observed in FNO on the linear regime.
  \item \textbf{Polynomial in $\eta$ by construction.} Each block
        expresses a learned $G_{n}$-shaped term; stacking blocks gives
        the architecture access to higher-order Craig--Sulem
        contributions --- the very terms that dominate when steepness
        or bandwidth is large.
\end{enumerate}
These together address the two strongest within-regime correlates of
rollout failure identified in \S\ref{sec:diagnostics}: bandwidth and
steepness.

\subsection{Implementation}
The base SpectralDNO uses one rank-$\mathrm{latent}$ spatial $\times$
spectral sandwich per block; the new architecture replaces it with
$n_{\mathrm{branches}}$ parallel \emph{pointwise} CS-shaped terms per
block. There is no tensor product between channel and branch axes ---
each branch carries one spatial channel and one $\xi$-multiplier, so
parameter count grows linearly in $n_{\mathrm{branches}}$. The
$\eta$-feature stack is shared across blocks; per-block parameters
are dominated by the two depth-aware multipliers and the
$\Phi$-projection.

\subsection{Cost and tunables}
A first run uses $\mathrm{width}=512$, $N_{\mathrm{blocks}}=8$,
$n_{\mathrm{branches}}=256$, $\mathrm{mult\_hidden}=128$ --- about
$1.1\times 10^{6}$ parameters, an order of magnitude smaller than the
FNO baseline (13.1M) and SpectralDNO (17.0M). The intent is to test
whether the inductive bias is doing the heavy lifting; capacity can be
scaled if the first signal is positive.

\subsection{Status}
Wired into the training and rollout-evaluation paths. Smoke checks
confirm: $y(2\xi) = 2\,y(\xi)$ to numerical precision (linearity in
$\xi$); $y(2\eta) - y(0) \ne 2(y(\eta) - y(0))$ after parameter
perturbation (nonlinearity in $\eta$); $y = G_{0}\xi$ exactly at
initialisation (zero-init blocks). 40-epoch training run in progress
(epoch 11 at the time of writing, best val $0.00199$ --- already 46\% below
the FNO baseline's all-time best of $0.00369$, attained in 1/7th the
training compute;
training is in parallel with the H-cons retrain on a separate GPU).

% ===================================================================
\section{Candidates for the next iteration}

These are not currently running but are next in the queue based on the
forensic findings and a literature scan.

\subsection{Multi-step unrolled training with gradients through all steps}
A direct extension of the one-step pushforward. Replace $k=1$
stop-gradient pushforward with a $k$-step unroll where gradients flow
through all steps, with a curriculum that ramps $k$ from $1$ to $8$:
\[
  u \leftarrow u, \quad \mathcal{L} \leftarrow 0,
  \qquad
  \text{for } j=1\dots k:\;
  u \leftarrow \Gop_\theta(u),\;
  \mathcal{L} \mathrel{+}= w_{j}\,\big\|u - u_{j}^{\mathrm{true}}\big\|^{2}.
\]
Stop-gradient one-step pushforward has a documented bias scaling like
$k^{2}$ relative to the true multi-step gradient, and $k\geq 8$ is often
needed for stable long rollouts. This is the most natural single change
to address the pushforward regression on random-sea regimes; memory
scales linearly in $k$, but at the current model size $k=4$ fits on a
single H100.

\subsection{Tendency-target reformulation}
Reframe the surrogate's output as a temporal derivative and integrate
with a real RK4 stepper at inference. For the DNO setting, the natural
place to apply this is on the $\xi$ side: instead of using the
closed-form Bernoulli RHS at rollout, have the surrogate also predict
$\partial_{t}\xi$ and stitch with RK4 externally. This decouples the
surrogate from carrying the dispersive phase rotation across one large
step.

\subsection{PDE-Refiner-style noise-conditioned refinement}
Standard MSE training under-weights low-amplitude high-$k$ modes ---
exactly the modes that drive modulational instability in Zakharov.
Replace one-shot prediction with $K$ iterative noise-conditioned
denoising passes (same model called $K$ times conditioning on a
log-noise channel). Inference cost goes up by $K\times$. The published
gain is reportedly smaller on FNO than U-Net, which is a concern for the
spectral backbone here but less so given the Craig--Sulem architecture
above.

\subsection{Rejected directions}
\begin{itemize}[leftmargin=1.3em,itemsep=2pt]
  \item Strict-symplectic / hard-Hamiltonian architectures --- wrong
        constraint level for this discretisation; the soft H-conservation
        loss already covers it.
  \item Global Lipschitz / spectral-norm constraints --- kills modulational
        dynamics for dispersive PDEs (true operator's $L^{2}$-induced
        norm is unitary).
  \item Simple capacity scaling --- current SpectralDNO is 17M params
        vs.\ FNO's 13M; capacity is not the bottleneck.
\end{itemize}

% ===================================================================
\section{Comparison of approaches}

\begin{center}
\footnotesize
\setlength{\tabcolsep}{4pt}
\begin{tabular}{@{}lllll@{}}
\toprule
Approach & Type & Constrains & Cost & Status \\
\midrule
Pushforward (1-step)       & loss   & rollout state dist.\          & $\sim 2\times$/step  & 40 ep done; 80 ep gated on GPU \\
H-cons (1-step)            & loss   & energy drift                  & $\sim 2\times$/step  & 25-ep retrain running \\
Craig--Sulem DNO           & model  & poly-$\eta$, lin-$\xi$        & $\sim 0.5\times$ base  & 40-ep run in progress (ep 11) \\
Multi-step unroll $k{=}4..8$ & loss & rollout state dist.\          & $k\times$/step       & queued \\
Tendency target + RK4      & model  & one-step ODE structure        & $4\times$ at rollout & queued \\
PDE-Refiner                & loss   & high-$k$ mode coverage        & $K\times$ at rollout & queued \\
\bottomrule
\end{tabular}
\end{center}

\noindent
The three running approaches are not mutually exclusive. The two
loss-side terms compose in the training script. Either can be paired
with the Craig--Sulem architecture without changes --- the model exposes
the same $\Gop_\theta(\eta,\xi,h)$ signature that the loss terms rely on.

% ===================================================================
\section{Open questions}

\begin{itemize}[leftmargin=1.3em,itemsep=3pt]
  \item Does the Craig--Sulem architecture without any rollout loss
        already beat pushforward-trained SpectralDNO? If so, the
        polynomial structure is the operative inductive bias, not the
        training loss.
  \item Does the random-sea regression under pushforward survive the
        Craig--Sulem architecture, or is it an artifact of the
        SpectralDNO single-sandwich design?
  \item Should the H-cons functional use $\Gop_{\mathrm{true}}\xi$ rather
        than $\Gop_{\theta}\xi$ when evaluating $H_{0}$? The current
        form rewards self-consistency of the surrogate's energy belief,
        not conservation of the true energy.
  \item Are the Benjamin--Feir g1 / modal NaN blow-ups (FNO baseline)
        driven by the $\eta$ update or the $\xi$ update? A step-by-step
        trace would say which loss term has more leverage.
\end{itemize}

% ===================================================================
\section*{Companion artifacts}

The per-regime rollout dashboard, including the error-vs-wavenumber
plots, is hosted at
\href{https://stunning-cocada-db3a7d.netlify.app}{stunning-cocada-db3a7d.netlify.app}.
Forensic reports (data audit, per-IC rollout breakdown, literature
review) and JSON statistics behind every claim in
\S\ref{sec:diagnostics} are available on request.

\end{document}
